% Figure: averaging drives the random component down as 1/sqrt(n)
% but leaves a systematic bias floor untouched. Two curves and a floor.
\begin{figure}[htbp]
\centering
\begin{tikzpicture}[
  axis/.style={draw=engblack, -{Stealth}, thick},
  randcurve/.style={draw=engblue, very thick},
  totalcurve/.style={draw=engred, very thick},
  floorline/.style={draw=engamber, thick, dashed},
  label/.style={font=\small},
]
% Axes
\draw[axis] (0,0) -- (11,0) node[anchor=west, label] {number of readings $n$};
\draw[axis] (0,0) -- (0,5.4) node[anchor=south, label] {reported uncertainty};

% x tick marks: n = 1, 4, 16, 64, 256 on a log-spaced axis
\foreach \x/\n in {0/1, 2.5/4, 5/16, 7.5/64, 10/256} {
  \draw (\x,-0.1) -- (\x,0.1);
  \node[font=\scriptsize, below] at (\x,-0.1) {\n};
}
% y tick marks
\foreach \y/\v in {0/0, 1/1, 2/2, 3/3, 4/4, 5/5} {
  \draw (-0.1,\y) -- (0.1,\y);
  \node[font=\scriptsize, left] at (-0.1,\y) {\v};
}

% Systematic-bias floor at u_B = 1.2 (constant)
\draw[floorline] (0,1.2) -- (10.5,1.2);
\node[engamber, font=\small, anchor=west] at (7.6,1.55)
  {systematic floor $u_B$};

% Random component u_A = 5 / sqrt(n). x maps n = 4^(x/2.5); plot the
% pure random 1/sqrt(n) decay through the five tick n-values.
% n:  1     4     16    64    256
% uA: 5.00  2.50  1.25  0.625 0.3125
\draw[randcurve]
  (0,5.0) .. controls (1.0,3.6) and (1.7,2.9) ..
  (2.5,2.5) .. controls (3.4,2.0) and (4.2,1.55) ..
  (5.0,1.25) .. controls (5.9,0.95) and (6.7,0.73) ..
  (7.5,0.625) .. controls (8.4,0.52) and (9.2,0.40) ..
  (10.0,0.3125);
\node[engblue, font=\small, anchor=west] at (5.3,0.55)
  {random $u_A = \sigma/\sqrt{n}$};

% Combined uncertainty u_c = sqrt(u_A^2 + u_B^2), asymptotes to floor.
% n:  1     4     16    64    256
% uc: 5.14  2.77  1.73  1.35  1.24
\draw[totalcurve]
  (0,5.14) .. controls (1.0,3.85) and (1.7,3.15) ..
  (2.5,2.77) .. controls (3.4,2.30) and (4.2,1.95) ..
  (5.0,1.73) .. controls (5.9,1.52) and (6.7,1.42) ..
  (7.5,1.35) .. controls (8.4,1.29) and (9.2,1.26) ..
  (10.0,1.24);
\node[engred, font=\small, anchor=west] at (2.7,3.15)
  {combined $u_c=\sqrt{u_A^2+u_B^2}$};

\end{tikzpicture}
\caption[Averaging floor: random decay against systematic bias]{Why
averaging reaches a floor. The random component of the reported
uncertainty falls as $\sigma/\sqrt{n}$ (blue) and would reach zero in
the limit of infinite readings. The systematic component $u_B$ (amber)
is a constant that averaging does not touch. The combined standard
uncertainty (red) is the quadrature sum $u_c=\sqrt{u_A^2+u_B^2}$; it
tracks the random curve while $u_A$ dominates, then flattens onto the
floor once $u_A$ falls below $u_B$. Past the point where the two curves
separate, further readings buy almost nothing, and the engineer who
keeps averaging is spending effort on the component that is already
small.}
\label{fig:vol01:ch01:averaging-floor}
\end{figure}
